\documentclass[10pt,a4paper]{amsart}
\usepackage[margin=19mm]{geometry}
\usepackage{amsmath,amssymb,mathtools}
\usepackage[scr=rsfs,cal=euler]{mathalfa}
\usepackage{xfrac}
\usepackage[unicode,hidelinks,bookmarks=false]{hyperref}
\newcommand{\ie}{i.e.\ }
\newcommand{\cf}{cf.\ }
\newcommand{\resp}{respectively\ }
\newcommand{\Subsection}{\subsection}
\providecommand{\nobreakdash}{\nobreak\hbox{-}}
\setlength{\emergencystretch}{2em}
\multlinegap0pt
\usepackage{bm}
\usepackage{bbm}
\usepackage{dsfont}
\usepackage{xspace}
\usepackage{cancel}
\usepackage{mathtools}
\usepackage{amsthm}
\makeatletter
\@ifundefined{theorem}{\newtheorem{theorem}{Theorem}}{}
\@ifundefined{lemma}{\newtheorem{lemma}[theorem]{Lemma}}{}
\makeatother
\newcommand{\ddt}[1]{\frac{\partial #1}{\partial t}}
\newcommand{\grad}{\nabla}
\newcommand{\mbf}[1]{\mathbf{#1}}
\title[An iterative approach to a fluid-rigid body interaction prob]{An iterative approach to a fluid-rigid body interaction problem}
\author{Charles M. Elliott, Thomas Sales}
\date{Source: arXiv:2601.13004v1 (CC BY 4.0)}
\begin{document}
\maketitle

\noindent\emph{Source and licence.} An iterative approach to a fluid-rigid body interaction problem. Version arXiv:2601.13004v1 (CC BY 4.0), distributed under the Creative Commons Attribution 4.0 International licence, as recorded in the arXiv licence metadata for this exact version and shown on the abstract page https://arxiv.org/abs/2601.13004v1. Full rights notice: licenses/arxiv-2601.13004-CC-BY-4.0.txt.\par\smallskip

\noindent\textit{Editorial scope.} Counted unit: the Lemma whose source label is \texttt{lemma: fsi energy estimates2} in the article (source0.tex lines 683--694, source label \texttt{lemma: fsi energy estimates2}), delivered together with the article's own complete original proof of it (source0.tex lines 695--720, one \texttt{proof} environment of 26 lines). No mathematics is written, repaired or re-derived: statement and proof are the article's own text line for line. Required context retained uncounted: eqn: leray L2 bound (lines 328--338); fsi energyestimate2 (lines 539--541). Every definition, display and label of those lines is retained unchanged. Omitted from this package and not redistributed: 1--327, 339--538, 542--682, 721--2437. Nothing inside a retained excerpt is omitted or altered: every retained body is delivered exactly as the article's lines are written, so no omission receipt of any kind accompanies this package. Citations: the retained excerpts carry no citation command at all, so no bibliography entry of the article is transcribed and no reference is redistributed. Reference adaptation: none was needed and none was made; every reference inside the delivered unit resolves inside it, so no unresolved reference or citation is produced. Printed slips kept literal: no printed word, symbol, hypothesis or number of the counted statement or of its proof was altered. Layout and class: the article's own class is \texttt{amsart} with packages including inputenc, amsfonts, amsmath, amsthm, graphicx, mathrsfs, verbatim, hyperref, array, float, appendix, caption, subcaption, tikz; the wrapper is the lane's pinned \texttt{amsart} editorial wrapper at 10pt on A4, which restates the theorem-like environments the retained text needs and adds the article's own macro definitions that the retained text needs (\texttt{\textbackslash ddt}, \texttt{\textbackslash grad}, \texttt{\textbackslash mbf}), with extra wrapper packages (bm, bbm, dsfont, xspace, cancel, mathtools). The delivered numbering is the unit's own. Assets: no retained span contains a figure environment, a graphics-inclusion command, a PDF-page inclusion, a table or a TikZ picture; no image file is copied into this package, the package declares no asset dependency and no image file is redistributed with it. Licence: version arXiv:2601.13004v1 of this article is distributed under the Creative Commons Attribution 4.0 International licence, as recorded in the arXiv licence metadata for this exact version and shown on the abstract page https://arxiv.org/abs/2601.13004v1. A full rights notice is delivered at licenses/arxiv-2601.13004-CC-BY-4.0.txt. Characters: every retained excerpt is pure ASCII, so the delivered engine, XeLaTeX with its default Latin Modern text font, reports no missing character.
\begin{lemma}  
    For $\mbf{u} \in \mbf{L}^2(\Omega(t))$ one has that
    \begin{align}
        \|P_\sigma \mbf{u}\|_{\mbf{L}^2(\Omega(t))} \leq \|\mbf{u}\|_{\mbf{L}^2(\Omega(t))}, \label{eqn: leray L2 bound}
    \end{align}
    and if one has further regularity, $\mbf{u} \in \mbf{H}^1(\Omega(t))$, then
    \begin{align}
        \|P_\sigma \mbf{u}\|_{\mbf{H}^1(\Omega(t))} \leq C\|\mbf{u}\|_{\mbf{H}^1(\Omega(t))}, \label{eqn: leray H1 bound}
    \end{align}
    for constants independent of $t \in [0,T]$.
\end{lemma}

	\begin{align}
		\frac{\rho}{2} \sup_{t \in [0,\widetilde{T}]} \int_{\Omega(t)} |\grad \overline{\mbf{u}}|^2 + \frac{\mu}{2} \int_0^{\widetilde{T}}\int_{\Omega(t)} |A_{\sigma} \overline{\mbf{u}}|^2 \leq  \mathcal{D}(\widetilde{T}), \label{fsi energyestimate2}
	\end{align}

\begin{lemma}
\label{lemma: fsi energy estimates2}
	On $[0,\widetilde{T}]$ for $\widetilde{T}$ as in the previous lemma, one has $\overline{\mbf{u}} \in H^1_{\mbf{L}^2}(\Phi)$, and there exists a unique pressure $\overline{p} \in L^2_{H^1 \cap L^2_0}(\Phi)$ such that
	\begin{equation}
        \begin{aligned}
		\int_{0}^{\widetilde{T}} \left(\left\| \ddt{\overline{\mbf{u}}} \right\|_{\mbf{L}^2(\Omega(t))}^2 + \|\overline{p}\|_{H^1(\Omega(t))}^2\right) &\leq \int_0^{\widetilde{T}} \|\mbf{B}\|_{\mbf{L}^2(\Omega(t))}^2\\
        &+C\left(1 + \int_0^{\widetilde{T}} \|\widetilde{\mbf{u}}\|_{\mbf{H}^2(\Omega(t))}^2 +  \sup_{t \in [0,\widetilde{T}]} \|\grad \widetilde{\mbf{u}}\|_{\mbf{L}^2(\Omega(t))}^2\right) \mathcal{D}(\widetilde{T})^2,
        \end{aligned}
        \label{fsi energyestimate3}
	\end{equation}
	for $\mathcal{D}(\cdot)$, $\mathcal{E}(\cdot)$ and $\kappa \in (0,1)$ as in the previous lemma.
\end{lemma}

\begin{proof}
	For almost all $t \in [0,\widetilde{T}]$ one finds that
    \[\rho P_\sigma \ddt{\overline{\mbf{u}}} = -\rho(\overline{\mbf{u}} \cdot \grad) \overline{\mbf{u}} - \mu A_{\sigma} \overline{\mbf{u}} - \rho (\widetilde{\mbf{u}} \cdot \grad)\overline{\mbf{u}} - \rho(\overline{\mbf{u}} \cdot \grad)\widetilde{\mbf{u}} - P_\sigma \mbf{B},\]
    and so
    \begin{multline*}
        \int_0^{\widetilde{T}} \left\| P_\sigma \ddt{\overline{\mbf{u}}} \right\|_{\mbf{L}^2(\Omega(t))}^2\\
        \leq C\int_0^{\widetilde{T}} \left(\|(\overline{\mbf{u}} \cdot \grad) \overline{\mbf{u}}\|_{\mbf{L}^2(\Omega(t))}^2 + \|A_{\sigma} \overline{\mbf{u}}\|_{\mbf{L}^2(\Omega(t))}^2 + \| (\widetilde{\mbf{u}} \cdot \grad)\overline{\mbf{u}}\|_{\mbf{L}^2(\Omega(t))}^2 + \| (\overline{\mbf{u}} \cdot \grad)\widetilde{\mbf{u}}\|_{\mbf{L}^2(\Omega(t))}^2 + \|P_\sigma \mbf{B}\|_{\mbf{L}^2(\Omega(t))}^2 \right),
    \end{multline*}
    where we note that
	\begin{align*}
		\int_0^T \|(\overline{\mbf{u}} \cdot \grad) \overline{\mbf{u}}\|_{\mbf{L}^2(\Omega(t))}^2 &\leq C\int_{0}^T \|\overline{\mbf{u}}\|_{\mbf{L}^\infty(\Omega(t))}^2 \|\grad \overline{\mbf{u}}\|_{\mbf{L}^2(\Omega(t))}^2\\
		&\leq C \left( \sup_{t \in [0,\widetilde{T}]} \|\grad \overline{\mbf{u}}\|_{\mbf{L}^2(\Omega(t))}^2 \right) \int_0^{\widetilde{T}} \|\overline{\mbf{u}}\|_{\mbf{H}^2(\Omega(t))}^2,
	\end{align*}
    where we have used the embedding $\mbf{H}^2(\Omega(t)) \hookrightarrow \mbf{L}^\infty(\Omega(t))$.
    We have similar estimates for $\int_0^{\widetilde{T}} \| (\widetilde{\mbf{u}} \cdot \grad)\overline{\mbf{u}}\|_{\mbf{L}^2(\Omega(t))}^2$ and $\int_0^{\widetilde{T}} \| (\overline{\mbf{u}} \cdot \grad)\widetilde{\mbf{u}}\|_{\mbf{L}^2(\Omega(t))}^2$.
	Thus by using \eqref{fsi energyestimate2} and \eqref{eqn: leray L2 bound} one finds that
    \begin{equation*}
        \int_0^{\widetilde{T}} \left\| P_\sigma \ddt{\overline{\mbf{u}}} \right\|_{\mbf{L}^2(\Omega(t))}^2 \leq \int_0^{\widetilde{T}} \|\mbf{B}\|_{\mbf{L}^2(\Omega(t))}^2 +C\left(1 + \int_0^{\widetilde{T}} \|\widetilde{\mbf{u}}\|_{\mbf{H}^2(\Omega(t))}^2 +  \sup_{t \in [0,\widetilde{T}]} \|\grad \widetilde{\mbf{u}}\|_{\mbf{L}^2(\Omega(t))}^2\right)\mathcal{D}(\widetilde{T})^2.
    \end{equation*}
	Likewise, for almost all $t \in [0,\widetilde{T}]$ one has that 
     \[\rho P_\sigma \ddt{\overline{\mbf{u}}} + \rho(\overline{\mbf{u}} \cdot \grad) \overline{\mbf{u}} - \mu A_{\sigma} \overline{\mbf{u}} + \rho (\widetilde{\mbf{u}} \cdot \grad)\overline{\mbf{u}} + \rho(\overline{\mbf{u}} \cdot \grad) \widetilde{\mbf{u}} - P_\sigma \mbf{B} \in \mbf{L}^2(\Omega(t)),\]
    where we use the Helmholtz decomposition to see this equals\footnote{Here we have written the non-solenoidal component as the pressure term, $\grad p$, and a sum of contributions from the terms involving $\overline{\mbf{u}}, \mbf{B}$.
    This is due to the fact that $P_\sigma$ does not commute with $\ddt{}$ or $\Delta$.}
    \[\grad \overline{p} + \rho\left(P_\sigma \ddt{\overline{\mbf{u}}} - \ddt{\overline{\mbf{u}}}\right) + (\mbf{B} - P_\sigma\mbf{B}) - \mu(\Delta \overline{\mbf{u}} + A_{\sigma}\overline{\mbf{u}}).\]
	The $L^2_{H^1}(\Phi)$ estimate on $\overline{p}$, and $L^2_{\mbf{L}^2}(\Phi)$ estimate for the non-solenoidal component of $\ddt{\overline{\mbf{u}}}$ now follow from using Poincar\'e's inequality and similar arguments to the above.
\end{proof}

\end{document}
